\documentclass[11pt,a4paper]{article}

\usepackage{iftex}
\RequirePDFTeX
\usepackage{cmap}
\usepackage[T2A]{fontenc}
\usepackage[utf8]{inputenc}
\usepackage[english,russian]{babel}
\usepackage{PTSerif}
\usepackage{amsmath,amssymb,microtype}
\usepackage[a4paper,left=23mm,right=23mm,top=20mm,bottom=21mm]{geometry}
\usepackage{xcolor,enumitem,fancyhdr,hyperref}

\definecolor{accent}{HTML}{187AA5}
\definecolor{ink}{HTML}{302B28}
\definecolor{muted}{HTML}{6D6560}
\definecolor{tint}{HTML}{EDF4F6}
\definecolor{rulecolor}{HTML}{A9CADA}
\hypersetup{unicode=true,colorlinks=true,urlcolor=accent,linkcolor=accent,
  pdftitle={Домашнее задание 2. Одноимпульсные манёвры},
  pdfauthor={Чиняев Владимир Николаевич}}
\urlstyle{same}
\color{ink}
\setlength{\parindent}{0pt}
\setlength{\parskip}{3pt}
\linespread{1.08}
\setlist[enumerate]{leftmargin=*,label=\textcolor{accent}{\arabic*.},
  itemsep=3pt,topsep=4pt,parsep=0pt}
\setlist[itemize]{leftmargin=*,itemsep=3pt,topsep=3pt,parsep=0pt}
\setlength{\emergencystretch}{1em}
\pagestyle{fancy}
\fancyhf{}
\renewcommand{\headrulewidth}{0pt}
\renewcommand{\footrulewidth}{0.4pt}
\renewcommand{\footrule}{\hbox to\headwidth{\color{rulecolor}\leaders\hrule height\footrulewidth\hfill}}
\fancyfoot[L]{\footnotesize\color{muted}Введение в маневрирование КА\quad /\quad ДЗ 2}
\fancyfoot[R]{\footnotesize\color{muted}\thepage}

\newcommand{\taskheading}[2]{%
  \par\vspace{5pt}{\large\bfseries\color{accent}Задание #1. #2}\par\smallskip}


\begin{document}

{\small\color{accent}ВВЕДЕНИЕ В МАНЕВРИРОВАНИЕ КА}\par
\vspace{3pt}
{\huge Домашнее задание 2}\par
{\large Одноимпульсные манёвры}\par
{\small\color{muted}К лекции 3}
\par{\small\bfseries\color{accent}Срок сдачи: 28 сентября 2026 года, 23:59}
\vspace{4pt}
{\color{rulecolor}\hrule}
\vspace{5pt}

Рассчитайте одноимпульсные переходы в центральном поле Земли.
Задания 1 и 2 обязательны, задание 3* --- по желанию (даёт дополнительные баллы).

\begingroup
\setlength{\fboxsep}{9pt}
\colorbox{tint}{\begin{minipage}{\dimexpr\linewidth-2\fboxsep\relax}
\textbf{Общие исходные данные}
\par\smallskip
{\centering
$\mu=398600{,}44158\ \text{км}^3/\text{с}^2,\qquad
R_\oplus=6378{,}137\ \text{км}.$\par\smallskip
$a_1=26000\ \text{км},\quad e_1=0{,}6,\quad
\omega_1=250^\circ,\quad i_1=64^\circ,\quad\Omega_1=0^\circ.$\par}
\smallskip
Сначала импульс прикладывается при $\nu_1=60^\circ$.
Рассматривайте некруговые эллипсы: $0<e_2<1$, $r_{\text{п},2}>R_\oplus$.
В пространственных задачах учитывайте и прямые, и обратные орбиты:
$0<i_2<180^\circ$.
\end{minipage}}
\endgroup

В коде используйте километры, секунды и радианы. В ответах углы $\omega,\Omega,\nu$
приводите к $[0,360^\circ)$, импульсы --- к м/с. Проверяйте $\mathbf r_2=\mathbf r_1$.

\taskheading{1}{Коррекция радиуса перицентра и аргумента перицентра}

Требуется получить $r_{\text{п},2}=10000$~км и $\omega_2=245^\circ$.
Плоскость орбиты и направление обращения сохраняются.

\begin{enumerate}[itemsep=2pt,topsep=2pt]
\item Найдите $u$ и $\nu_2$. Из $r_{\text{п},2}=a_2(1-e_2)$ и уравнения орбиты
выведите $e_2,a_2$. Укажите условия существования допустимого решения.
\item При $\nu_1=60^\circ$ найдите конечные элементы, вектор импульса в
инерциальных осях и его модуль. Проверьте общую точку и целевые элементы.
\item Переберите $0\le\nu_1<360^\circ$ с шагом не более $0{,}5^\circ$.
Постройте $\Delta v(\nu_1)$ с разрывами там, где допустимого эллипса нет.
Найдите наименьшие затраты на сетке; при желании уточните минимум.
\item Сравните найденный минимум с затратами при $60^\circ$.
Объясните, почему этот набор элементов можно независимо скорректировать
не в любой точке орбиты.
\end{enumerate}

\taskheading{2}{Коррекция большой полуоси, эксцентриситета\newline и аргумента перицентра}

При $\nu_1=60^\circ$ требуется получить
$a_2=28000$~км, $e_2=0{,}65$, $\omega_2=240^\circ$.
Наклонение и ДВУ заранее не задаются.

\begin{enumerate}[itemsep=2pt,topsep=2pt]
\item Найдите возможные $\nu_2$ из уравнения радиуса.
Для каждой определите $u_2=\omega_2+\nu_2$ и допустимые $i_2$ из
$r_z=r\sin i_2\sin u_2$. Затем найдите $\Omega_2$ по первым двум
координатам радиус-вектора. Учитывайте четверть угла и все допустимые ветви.
\item Найдите все конечные орбиты и импульсы, включая варианты с $i_2>90^\circ$.
Составьте таблицу элементов и затрат при $\nu_1=60^\circ$.
Проверьте заданные элементы, общую точку и радиус перицентра.
\item Переберите $0\le\nu_1<360^\circ$ с шагом не более $0{,}5^\circ$.
Постройте на одном графике $\Delta v(\nu_1)$ для прямой и обратной ветвей.
Недопустимые точки пропускайте. Найдите наиболее экономичный манёвр.
\item Покажите исходную и найденные орбиты на одном пространственном рисунке.
Объясните, как изменение плоскости позволяет независимо задать $a,e,\omega$.
\end{enumerate}

\newpage
{\small\color{muted}Домашнее задание 2\hfill К лекции 3}

\taskheading{3*}{Коррекция большой полуоси, эксцентриситета и наклонения}
\textit{Необязательное задание.}

Требуется получить
\[
a_2=28000\ \text{км},\qquad e_2=0{,}65,\qquad i_2=70^\circ.
\]
ДВУ $\Omega_2$ и аргумент перицентра $\omega_2$ определяются из решения.
При заданном наклонении через точку манёвра могут проходить две разные
орбитальные плоскости. В каждой из них заданные $a_2,e_2$ допускают
два положения перицентра. Для выбранных данных получаются четыре различных перехода.

\begin{enumerate}
\item При $\nu_1=60^\circ$ решите условие
$\widehat{\mathbf h}_2\cdot\mathbf r=0$ относительно $\Omega_2$.
Используйте выражение
\[
\widehat{\mathbf h}_2=
(\sin i_2\sin\Omega_2,\;-\sin i_2\cos\Omega_2,\;\cos i_2)^{\mathsf T}.
\]
В каждой найденной плоскости определите обе конечные орбиты с заданными
$a_2,e_2$. При расчёте $\omega_2$ заново найдите аргумент широты в этой плоскости.
\item Обозначьте решения $1A,1B,2A,2B$: номер соответствует одной из двух
плоскостей, буква --- одной из двух конечных истинных аномалий.
Представьте таблицу элементов и импульсов при $60^\circ$.
Проверьте общую точку, целевые элементы и допустимость всех четырёх орбит.
\item Переберите полный оборот $0\le\nu_1<360^\circ$ с шагом не более
$0{,}5^\circ$. Постройте на одном графике четыре линии $\Delta v(\nu_1)$
с легендой. Сохраняйте происхождение ветвей при переходе углов через
$0/360^\circ$: не переупорядочивайте их по текущей ДВУ или величине импульса.
\item Найдите наилучший результат для каждой ветви и среди всех четырёх.
\item Объясните происхождение четырёх решений. При каких условиях две
плоскости или две орбиты в одной плоскости сливаются?
Почему для произвольных исходных и целевых данных число решений может быть меньше четырёх?
\end{enumerate}

\vspace{4pt}
\textbf{Работа с ноутбуком}

Файл \texttt{hw\_2.ipynb} содержит полные решения численных примеров лекции,
функции решения тригонометрического и квадратного уравнений, пересчёта
орбитальных элементов в векторы состояния, построения рисунков и графиков.
В конце файла подготовлены ячейки для ваших решений.

Все вспомогательные функции уже написаны в ноутбуке. Вам нужно корректно
объединить их вызовы, задать условия отбора решений и объяснить полученные результаты.
Для поиска места манёвра сначала переберите всю орбиту и сравните все ветви.
Локальное уточнение выполняйте после перебора.

\vspace{5pt}
{\color{rulecolor}\hrule}
\vspace{5pt}
\textbf{Что представить}

Сдайте один ноутбук с решёнными заданиями 1 и 2 и, по желанию, заданием~3*.
Формулы, рассуждения и ответы оформите в текстовых ячейках, расчёты и графики
--- в кодовых. Перед отправкой перезапустите ядро, выполните все ячейки
по порядку и сохраните ноутбук с результатами вычислений и графиками.

Название файла для отправки: \texttt{hw\_2\_name\_surname.ipynb}.
Вместо \texttt{name\_surname} укажите свои имя и фамилию латиницей.

\end{document}
